\chapter{Background}
\label{ch:background}

This chapter introduces the theoretical and methodological background required
for the development of our guidance approach.

We begin with flow matching, the generative framework underlying ArchesWeather,
and introduce the notion of probability paths, noisy states, velocity fields, and clean
estimates used throughout this work.
We then formulate guidance from a posterior-sampling perspective and derive how
conditioning modifies the flow-matching velocity field.

Next, we describe the components of ArchesWeather relevant to our method,
including its deterministic forecast, residual flow model, and sampling
procedure.
Finally, we review existing work on guidance for generative weather models and
position our approach within the current literature.


\section{Flow matching}

Generative modeling frames the task of generating objects
\(x_1 \in \mathbb{R}^d\) as sampling from a data distribution \(p_1\).
Flow matching \parencite{lipman2023flow} models \(p_1\) as the endpoint of a
continuous probability path
\[
    (p_t)_{t\in[0,1]},
\]
which transports a tractable source distribution \(p_0\), typically a standard
Gaussian, to the data distribution \(p_1\).

At flow time \(t\), we denote the current state along this path by \(z_t\),
with
\[
    z_0 \sim p_0,
    \qquad
    z_t \sim p_t,
    \qquad
    z_1 = x_1 \sim p_1.
\]
Throughout this work, we refer to $z_0$ as the \emph{noise}, intermediate states
$z_t$, $t\in(0,1)$, as \emph{noisy states}, and the endpoint $z_1=x_1$ as the
\emph{clean state}.

The evolution of \(z_t\) is governed by a time-dependent velocity field
\(u_t\), defining the ordinary differential equation
\begin{equation}
    \frac{\mathrm{d}}{\mathrm{d}t} z_t
    =
    u_t(z_t).
    \label{eq:flow-ode}
\end{equation}
Starting from \(z_0\sim p_0\) and integrating this ODE from \(t=0\) to \(t=1\)
therefore transforms samples from the source distribution into samples from
the data distribution.

% Equivalently, the velocity field and probability path satisfy the continuity
% equation
% \begin{equation}
%     \partial_t p_t(z)
%     +
%     \nabla_z \cdot
%     \bigl(
%         p_t(z)u_t(z)
%     \bigr)
%     =
%     0,
%     \label{eq:continuity-equation}
% \end{equation}
% which describes how probability mass is transported along the flow.


\subsection{Conditional probability paths}

Directly specifying the marginal probability path \(p_t\) and its corresponding
velocity field \(u_t\) is generally intractable.
Flow matching circumvents this problem by defining a tractable
\emph{conditional probability path}
\[
    p_t(\cdot\mid x_1)
\]
for each data sample \(x_1\sim p_1\).

A common choice is the Gaussian conditional probability path
\begin{equation}
    p_t(\cdot\mid x_1)
    =
    \mathcal{N}
    \left(
        \alpha_t x_1,\,
        \beta_t^2 I
    \right),
    \label{eq:gaussian-probability-path}
\end{equation}
where the schedules \(\alpha_t,\beta_t\in[0,1]\) interpolate between noise and
data according to
\[
    \alpha_0=0,
    \qquad
    \beta_0=1,
    \qquad
    \alpha_1=1,
    \qquad
    \beta_1=0.
\]
Consequently,
\begin{align}
    p_0(\cdot\mid x_1)
    &= \mathcal{N}(0,I),
    \\
    p_1(\cdot\mid x_1)
    &= \delta_{x_1}.
\end{align}

Sampling from the conditional path is equivalent to interpolating a data
sample \(x_1\sim p_1\) with independent Gaussian noise
\(\varepsilon\sim\mathcal{N}(0,I)\):
\begin{equation}
    z_t
    =
    \alpha_t x_1
    +
    \beta_t \varepsilon
    \sim
    p_t(\cdot\mid x_1).
    \label{eq:gaussian-interpolation}
\end{equation}

Marginalizing over the data distribution gives the corresponding marginal
probability path
\begin{equation}
    p_t(z_t)
    =
    \int
    p_t(z_t\mid x_1)\,
    p_1(x_1)\,
    \mathrm{d}x_1.
    \label{eq:marginal-probability-path}
\end{equation}


\subsection{Conditional and marginal velocity fields}

Each conditional probability path \(p_t(\cdot\mid x_1)\) is generated by a
corresponding conditional velocity field \(u_t(z_t\mid x_1)\).

For the Gaussian path in \Cref{eq:gaussian-probability-path}, differentiating
the interpolation in \Cref{eq:gaussian-interpolation} gives
\begin{equation}
    u_t(z_t\mid x_1)
    =
    \dot{\alpha}_t x_1
    +
    \dot{\beta}_t \varepsilon.
    \label{eq:conditional-velocity-noise}
\end{equation}
Substituting
\[
    \varepsilon
    =
    \frac{z_t-\alpha_t x_1}{\beta_t}
\]
yields
\begin{equation}
    u_t(z_t\mid x_1)
    =
    \frac{\dot{\beta}_t}{\beta_t}z_t
    +
    \left(
        \dot{\alpha}_t
        -
        \frac{\dot{\beta}_t}{\beta_t}\alpha_t
    \right)x_1.
    \label{eq:conditional-velocity}
\end{equation}

The marginal velocity field generating \(p_t\) can then be expressed as the
conditional velocity averaged over all clean samples \(x_1\) compatible with
the current noisy state \(z_t\):
\begin{equation}
    u_t(z_t)
    =
    \mathbb{E}_{x_1\sim p_t(\cdot\mid z_t)}
    \left[
        u_t(z_t\mid x_1)
    \right].
    \label{eq:marginal-velocity}
\end{equation}

While this marginal velocity is generally not available in closed form, the
conditional velocity in \Cref{eq:conditional-velocity} is tractable during
training, since both \(x_1\) and the sampled noise \(\varepsilon\) are known.


\subsection{Flow-matching objective}

Flow matching exploits the fact that a neural network can be trained on the
tractable conditional velocity field while learning the corresponding
marginal velocity field \parencite{lipman2023flow}.
Writing \(u_\theta(z_t,t)\) for the learned velocity field, the conditional
flow-matching objective is
\begin{equation}
    \mathcal{L}_{\mathrm{FM}}(\theta)
    =
    \mathbb{E}_{
        \substack{
            t\sim\mathcal{U}[0,1],\\
            x_1\sim p_1,\\
            z_t\sim p_t(\cdot\mid x_1)
        }
    }
    \left[
        \left\|
            u_\theta(z_t,t)
            -
            u_t(z_t\mid x_1)
        \right\|_2^2
    \right].
    \label{eq:fm-objective}
\end{equation}
This conditional objective has the same gradient with respect to the model
parameters as the corresponding, but intractable, marginal flow-matching
objective \parencite{lipman2023flow}.

After training, \(u_\theta(z_t,t)\) approximates the marginal velocity field
\(u_t(z_t)\).
Generation then consists of sampling
\[
    z_0\sim p_0
\]
and numerically integrating
\begin{equation}
    \frac{\mathrm{d}}{\mathrm{d}t}z_t
    =
    u_\theta(z_t,t)
\end{equation}
from \(t=0\) to \(t=1\), such that the resulting endpoint
\(z_1\) approximately follows \(p_1\).


\subsection{Linear probability path}

A particularly simple choice is the linear, or optimal-transport, probability
path
\begin{equation}
    z_t
    =
    t x_1
    +
    (1-t)\varepsilon,
    \qquad
    \varepsilon\sim\mathcal{N}(0,I),
    \label{eq:linear-probability-path}
\end{equation}
corresponding to
\[
    \alpha_t=t,
    \qquad
    \beta_t=1-t.
\]
Its conditional velocity reduces to the constant displacement
\begin{equation}
    u_t(z_t\mid x_1)
    =
    x_1-\varepsilon.
    \label{eq:linear-conditional-velocity}
\end{equation}
Training therefore amounts to least-squares regression onto the displacement
between the source noise and the clean sample.

At an intermediate flow time \(t<1\), the noisy state \(z_t\) does not uniquely
identify the clean sample \(x_1\) from which it originates.
Instead, the possible clean samples are described by the conditional
distribution
\[
    p_t(x_1\mid z_t).
\]
In practice, a generative model can provide a tractable clean estimate
\begin{equation}
    \hat{x}_t
    :=
    \hat{x}_\theta(z_t,t)
    \label{eq:clean-estimate}
\end{equation}
from the current noisy state.
The distinction between the full conditional distribution
\(p_t(x_1\mid z_t)\) and the point estimate \(\hat{x}_t\) will become important
when introducing guidance in the following section.

\section{Guidance}

In unconditional generation, the flow-matching model produces samples from the
data distribution \(p_1(x_1)\).
In many applications, however, we are interested in generating samples that
satisfy a desired condition \(y\).
Guidance modifies the sampling process to favor samples from the corresponding
conditional distribution
\[
    p_1(x_1\mid y).
\]

We first formulate guidance from a posterior-sampling perspective and then
translate the resulting score modification to the flow-matching velocity
field.
Finally, we introduce a general guidance-strength schedule that will be used
throughout this work.


\subsection{Guidance as posterior sampling}

Rather than considering only the conditional distribution at the clean endpoint,
we can define the corresponding distribution over noisy states at any flow time
\(t\).
By Bayes' rule,
\begin{equation}
    p_t(z_t\mid y)
    \propto
    p_t(z_t)\,p_t(y\mid z_t).
\end{equation}
Taking the gradient with respect to the noisy state gives the conditional score
\begin{align}
    \nabla_{z_t}\log p_t(z_t\mid y)
    &=
    \nabla_{z_t}\log p_t(z_t)
    +
    \nabla_{z_t}\log p_t(y\mid z_t).
    \label{eq:conditional-score}
\end{align}

The first term corresponds to the unconditional generative process, while the
second introduces the desired condition \(y\).
Guidance can therefore be interpreted as modifying the unconditional
generation process through the likelihood gradient
\[
    \nabla_{z_t}\log p_t(y\mid z_t).
\]

In general, however, the intermediate likelihood \(p_t(y\mid z_t)\) is not
directly available.
We therefore marginalize over the clean sample \(x_1\) and assume that the
condition \(y\) depends on the generative process only through \(x_1\):
\begin{align}
    p_t(y\mid z_t)
    &=
    \int
    p_t(y\mid x_1,z_t)\,
    p_t(x_1\mid z_t)\,
    \mathrm{d}x_1
    \\
    &=
    \int
    p(y\mid x_1)\,
    p_t(x_1\mid z_t)\,
    \mathrm{d}x_1,
    \qquad
    y \perp z_t \mid x_1
    \\
    &=
    \mathbb{E}_{x_1\sim p_t(\cdot\mid z_t)}
    \left[
        p(y\mid x_1)
    \right]
    \\
    &\approx
    \int
    p(y\mid x_1)\,
    \delta_{\hat{x}_t}(x_1)\,
    \mathrm{d}x_1,
    \qquad
    p_t(x_1\mid z_t)
    \approx
    \delta_{\hat{x}_t}(x_1)
    \\
    &=
    p(y\mid\hat{x}_t).
    \label{eq:clean-likelihood-approximation}
\end{align}

Here, \(\hat{x}_t=\hat{x}_\theta(z_t,t)\) denotes the clean estimate obtained
from the current noisy state.
The approximation replaces the full conditional distribution
\(p_t(x_1\mid z_t)\) by a point mass at \(\hat{x}_t\), yielding
\begin{equation}
    p_t(y\mid z_t)
    \approx
    p(y\mid\hat{x}_t).
    \label{eq:likelihood-approximation}
\end{equation}

The likelihood \(p(y\mid x_1)\) can in principle be modeled in different ways.
A convenient choice for gradient-based guidance is to define it through a
differentiable loss \(\mathcal{L}\):
\begin{equation}
    p(y\mid x_1)
    \propto
    \exp\left(
        -\mathcal{L}(x_1,y)
    \right).
    \label{eq:energy-likelihood}
\end{equation}
It follows that
\begin{equation}
    \nabla_{z_t}\log p(y\mid\hat{x}_t)
    =
    -
    \nabla_{z_t}\mathcal{L}(\hat{x}_t,y),
    \label{eq:likelihood-loss}
\end{equation}
where the gradient is propagated through the clean estimate
\(\hat{x}_t=\hat{x}_\theta(z_t,t)\).


\subsection{From scores to flow-matching velocities}

The previous derivation expresses guidance in terms of scores.
To apply it to flow matching, we require the corresponding modification of the
velocity field.

For the Gaussian probability path introduced in
\Cref{eq:gaussian-probability-path}, the marginal velocity field and score are
related by
\begin{equation}
    u_t(z_t)
    =
    a_t
    \nabla_{z_t}\log p_t(z_t)
    +
    b_t z_t,
    \label{eq:score-velocity-relation}
\end{equation}
where
\begin{align}
    a_t
    &:=
    \beta_t^2
    \left(
        \frac{\dot{\alpha}_t}{\alpha_t}
        -
        \frac{\dot{\beta}_t}{\beta_t}
    \right),
    \\
    b_t
    &:=
    \frac{\dot{\alpha}_t}{\alpha_t}.
\end{align}

Applying the same relation to the conditional distribution gives
\begin{equation}
    u_t(z_t\mid y)
    =
    a_t
    \nabla_{z_t}\log p_t(z_t\mid y)
    +
    b_t z_t.
\end{equation}
Substituting the conditional-score decomposition yields
\begin{align}
    u_t(z_t\mid y)
    &=
    a_t
    \left[
        \nabla_{z_t}\log p_t(z_t)
        +
        \nabla_{z_t}\log p_t(y\mid z_t)
    \right]
    +
    b_t z_t
    \\
    &=
    u_t(z_t)
    +
    a_t
    \nabla_{z_t}\log p_t(y\mid z_t).
    \label{eq:conditional-velocity}
\end{align}

Using the clean-estimate approximation derived above, we obtain
\begin{equation}
    u_t(z_t\mid y)
    \approx
    u_t(z_t)
    +
    a_t
    \nabla_{z_t}\log p(y\mid\hat{x}_t).
    \label{eq:conditional-velocity-approximation}
\end{equation}

Under the loss-based likelihood in \Cref{eq:energy-likelihood}, this can
equivalently be written as
\begin{equation}
    u_t(z_t\mid y)
    \approx
    u_t(z_t)
    -
    a_t
    \nabla_{z_t}\mathcal{L}(\hat{x}_t,y).
    \label{eq:conditional-velocity-loss}
\end{equation}

Thus, gradient-based guidance in flow matching consists of evaluating a loss
on the clean estimate \(\hat{x}_t\), differentiating it with respect to the
current noisy state \(z_t\), and using the resulting gradient to modify the
velocity field.


\subsection{Guidance-strength schedules}

The coefficient \(a_t\) in \Cref{eq:conditional-velocity-loss} is determined
by the underlying probability path.
Using this coefficient preserves the posterior-sampling interpretation, up to
the clean-estimate approximation introduced above.

In practice, guidance methods often rescale this contribution to obtain
stronger or weaker control over the generated samples.
We therefore replace \(a_t\) by a general guidance-strength schedule
\begin{equation}
    \lambda_t
    :=
    w\,\gamma_t\,c_t,
    \label{eq:guidance-strength-schedule}
\end{equation}
where
\begin{itemize}[noitemsep]
    \item \(w \geq 0\) controls the overall guidance strength;
    \item \(\gamma_t \geq 0\) controls how guidance varies along the flow;
    \item \(c_t > 0\) is an optional algorithm-specific normalization factor.
\end{itemize}

We then define the \emph{guided velocity field} as
\begin{align}
    \tilde{u}_t(z_t\mid y)
    &:=
    u_t(z_t)
    +
    \lambda_t
    \nabla_{z_t}\log p(y\mid\hat{x}_t)
    \\
    &=
    u_t(z_t)
    -
    \lambda_t
    \nabla_{z_t}\mathcal{L}(\hat{x}_t,y).
    \label{eq:guided-velocity}
\end{align}
We refer to \(u_t(z_t)\) as the \emph{unguided velocity field}.

Setting
\[
    \lambda_t=a_t
\]
recovers the approximate conditional velocity field in
\Cref{eq:conditional-velocity-approximation}.
For a general choice of \(\lambda_t\), however, the resulting dynamics no
longer necessarily correspond to sampling from \(p_1(x_1\mid y)\).
Instead, the guidance strength becomes an algorithmic parameter controlling
the trade-off between the original generative dynamics and the imposed
condition.

\section{ArchesWeather}

ArchesWeather \parencite{couairon2024archesweather} is a probabilistic
medium-range weather forecasting model that combines a deterministic forecast
with a flow-matching model of the forecast residual.
In this section, we describe the components required for the development
of our guidance method and refer to the original work for architectural and
training details.


\subsection{Dataset}

ArchesWeather is trained on the ERA5 reanalysis dataset
\parencite{hersbach2020era5}, using one of the versions preprocessed by
WeatherBench~2 \parencite{rasp2024weatherbench2}.
The data are represented on a regular \(1.5^\circ\) latitude--longitude grid
with \(121\times240\) spatial points.

Each weather state
\[
    x_n \in \mathbb{R}^{82\times121\times240}
\]
contains \(82\) atmospheric fields.
Four are surface variables: \(2\)~m temperature, mean sea-level pressure, and
the two components (\(u\) and \(v\)) of the \(10\)~m wind.
The remaining \(78\) fields correspond to six upper-air variables
\[
    \{
        \text{geopotential},
        \text{temperature},
        \text{specific humidity},
        u,
        v,
        \text{vertical velocity}
    \}
\]
evaluated at the \(13\) pressure levels
\[
    \{50,100,150,200,250,300,400,500,600,700,850,925,1000\}
    \ \mathrm{hPa}.
\]

Weather states are available at \(00{:}00\), \(06{:}00\), \(12{:}00\), and
\(18{:}00\) UTC.
Following the WeatherBench~2 evaluation protocol, dates from January 2020
onward are reserved for evaluation.

Both the deterministic and generative components operate in normalized data
space.
Unless stated otherwise, the variables introduced below therefore refer to
their normalized representation.


\subsection{Forecasting setup}

ArchesWeather generates forecast trajectories autoregressively.
We index forecast lead time by \(n\), which we refer to as
\emph{weather time}.
At each weather step, the model produces a forecast for the next state
\(x_{n+1}\).

Rather than generating \(x_{n+1}\) directly, ArchesWeather decomposes the
forecast into a deterministic component and a stochastic residual component.
The deterministic model first produces a base forecast
\[
    \hat{x}^{\mathrm{det}}_{n+1}.
\]
The corresponding ground-truth residual is defined as
\begin{equation}
    r_{n+1}
    :=
    x_{n+1}
    -
    \hat{x}^{\mathrm{det}}_{n+1}.
    \label{eq:arches-residual}
\end{equation}

A conditional flow-matching model is then trained to generate this residual.
At inference time, a sampled residual \(\hat r_{n+1}\) is combined with the
deterministic prediction as
\begin{equation}
    \hat{x}_{n+1}
    =
    \hat{x}^{\mathrm{det}}_{n+1}
    +
    \hat{r}_{n+1}.
    \label{eq:arches-forecast}
\end{equation}

Sampling different residuals while keeping the deterministic forecast fixed
therefore produces different members of an ensemble forecast.


\subsection{Residual flow model}

Generating each stochastic residual requires a complete flow-matching
trajectory.
ArchesWeather therefore involves two nested notions of time: the weather-time
index \(n\), which advances the autoregressive forecast, and the flow time
\(t\), which describes the generation of a residual within each weather step.

Following the notation introduced in the previous sections, the clean endpoint
of the flow is now the forecast residual \(r_{n+1}\), while \(z_t\) denotes its
noisy state at flow time \(t\).

ArchesWeather parameterizes the Gaussian probability path in terms of a noise
level \(s_t\), which decreases linearly from \(1\) at the beginning of the flow
to \(0\) at its endpoint.
The noisy residual is constructed as
\begin{equation}
    z_t
    =
    (1-s_t)\,r_{n+1}
    +
    s_t\,\varepsilon,
    \qquad
    \varepsilon\sim\mathcal{N}(0,I).
    \label{eq:arches-probability-path}
\end{equation}

This corresponds directly to the Gaussian probability path introduced in
\Cref{eq:gaussian-probability-path}, with
\begin{equation}
    \alpha_t = 1-s_t,
    \qquad
    \beta_t = s_t.
    \label{eq:arches-path-mapping}
\end{equation}
The flow therefore starts from Gaussian noise and progressively approaches the
clean residual \(r_{n+1}\).

Unlike the generic flow-matching formulation, in which the neural network
directly predicts the velocity field, the ArchesWeather residual model is
trained to predict the clean residual from the noisy state.
We denote this estimate by
\begin{equation}
    \hat{r}_t
    :=
    r_\theta(z_t,t,c_n),
    \label{eq:arches-clean-residual}
\end{equation}
where
\[
    c_n
    =
    \left(
        x_{n-1},
        x_n,
        \hat{x}^{\mathrm{det}}_{n+1}
    \right)
\]
collects the conditioning information provided to the residual model.

The model is trained using the clean-prediction objective
\begin{equation}
    \mathcal{L}_{\mathrm{Arches}}(\theta)
    =
    \mathbb{E}_{t,r_{n+1},\varepsilon}
    \left[
        \left\|
            r_\theta(z_t,t,c_n)
            -
            r_{n+1}
        \right\|_2^2
    \right].
    \label{eq:arches-training-objective}
\end{equation}

In the notation of the previous sections, \(\hat r_t\) thus plays the role of
the clean estimate \(\hat x_t\).
Since the flow generates a residual rather than the complete weather state,
the corresponding clean weather estimate at flow time \(t\) is
\begin{equation}
    \hat{x}_t
    :=
    \hat{x}^{\mathrm{det}}_{n+1}
    +
    \hat{r}_t.
    \label{eq:arches-clean-state}
\end{equation}
This quantity will later provide the state on which our guidance objectives
are evaluated.


\subsection{Sampling}

Although the network predicts the clean residual, sampling requires the
corresponding velocity field.
From \Cref{eq:arches-probability-path}, the noise associated with the clean
estimate \(\hat r_t\) is
\begin{equation}
    \hat{\varepsilon}_t
    =
    \frac{
        z_t-(1-s_t)\hat r_t
    }{s_t}.
    \label{eq:arches-noise-estimate}
\end{equation}

For the linear path
\[
    \alpha_t=1-s_t,
    \qquad
    \beta_t=s_t,
\]
the velocity is the displacement between the clean residual estimate and the
corresponding noise estimate:
\begin{align}
    u_\theta(z_t,t,c_n)
    &=
    \hat r_t-\hat{\varepsilon}_t
    \\
    &=
    \hat r_t
    -
    \frac{
        z_t-(1-s_t)\hat r_t
    }{s_t}
    \\
    &=
    \frac{
        \hat r_t-z_t
    }{s_t}.
    \label{eq:arches-velocity}
\end{align}

ArchesWeather discretizes the flow into \(T=25\) sampling steps.
In the following, \(t=0,\ldots,T\) denotes the corresponding discrete
flow-step index.
Starting from Gaussian noise
\[
    z_0 \sim \mathcal{N}(0,I),
\]
the residual is generated through Euler integration,
\begin{equation}
    z_{t+1}
    =
    z_t
    +
    h_t\,u_\theta(z_t,t,c_n),
    \qquad
    h_t=s_t-s_{t+1},
    \qquad
    t=0,\ldots,T-1.
    \label{eq:arches-euler}
\end{equation}
Since the noise level decreases from \(s_0=1\) to \(s_T=0\), the sequence
\(\{z_t\}_{t=0}^{T}\) progressively evolves from noise toward a clean
residual.

After the final flow step, the endpoint is taken as the sampled residual,
\begin{equation}
    \hat{r}_{n+1}
    :=
    z_T,
    \label{eq:arches-sampled-residual}
\end{equation}
which is added to the deterministic forecast to obtain
\begin{equation}
    \hat{x}_{n+1}
    =
    \hat{x}^{\mathrm{det}}_{n+1}
    +
    \hat{r}_{n+1}
    =
    \hat{x}^{\mathrm{det}}_{n+1}
    +
    z_T.
    \label{eq:arches-final-forecast}
\end{equation}

The resulting nested time structure is illustrated in \Cref{fig:time-scales},
while \Cref{fig:arches-computational-graph} summarizes the computations
performed within a single weather step.

\begin{figure}[H]
    \centering
    \resizebox{\textwidth}{!}{
        \begin{tikzpicture}[x=0.9cm,y=0.9cm]
            \def\L{3.4}
            \def\M{25}

            \draw (-0.6,0) -- (0,0);
            \draw[-{Latex[length=3mm]}] (0,0) -- (10.8,0);

            \foreach \j in {0,1,2} {
                \foreach \i in {1,...,24} {
                    \draw[line width=0.3pt]
                    ({\j*\L + \i*\L/\M}, -0.10)
                    --
                    ({\j*\L + \i*\L/\M}, 0.10);
                }
            }

            \foreach \x/\lab in {
                0/{n=0},
                3.4/{n=1},
                6.8/{\cdots},
                10.2/{n=N}
            }{
                \draw[line width=1pt] (\x,-0.28) -- (\x,0.28);
                \node[below=6pt] at (\x,-0.28)
                {\scriptsize $\lab$};
            }

            \node[above] at (1.7,0.3)
            {\scriptsize $T$ flow steps};
        \end{tikzpicture}
    }
    \caption{
        The two nested time scales in ArchesWeather.
        Each autoregressive weather step \(n\rightarrow n+1\) contains a
        complete flow of \(T=25\) steps used to generate the forecast residual.
    }
    \label{fig:time-scales}
\end{figure}

\begin{figure}[H]
    \centering
    \resizebox{\textwidth}{!}{
        % Computational graph: deterministic core + generative residual flow.
% Self-contained tikzpicture (styles inline); requires the kugreen* colors and
% \usetikzlibrary{arrows.meta, positioning, calc}. Shared by presentation.tex
% (Model frame) and standalone.tex (PNG export).
\begin{tikzpicture}[
  font=\small,
  state/.style={draw=kugreen, fill=kugreenlyslyslys, rounded corners=2pt,
                minimum height=0.72cm, inner sep=4pt},
  model/.style={draw=kugreen, fill=kugreen, text=white, rounded corners=2pt,
                minimum height=0.9cm, inner sep=5pt},
  emb/.style={draw=kugreenlys, fill=kugreenlyslyslys!40, rounded corners=2pt,
              inner sep=3pt, font=\scriptsize},
  op/.style={draw=kugreen, circle, fill=white, inner sep=1pt},
  flow/.style={-{Latex[length=2.2mm]}, kugreen!80!black, thick},
  loop/.style={-{Latex[length=2.2mm]}, kugreen!80!black, thick, dashed},
  lbl/.style={font=\scriptsize, text=kugreen!60!black},
]

% ===== lane 1: deterministic pipeline (inputs on the same row) =====
\node[state] (xprev) at (0,3.0) {$x_{n-1}$};
\node[state] (xcur)  at (1.1,3.0) {$x_{n}$};
\node[op]    (catdet) at (2.05,3.0) {$\|$};
\node[model] (det)   at (3.9,3.0) {deterministic model};
\node[state] (xdet)  at (6.6,3.0) {$\hat{x}^{\text{det}}$};

\draw[flow] (xprev.north) to[out=55, in=125] (catdet.north);
\draw[flow] (xcur) -- (catdet);
\draw[flow] (catdet) -- (det);
\draw[flow] (det) -- (xdet);

% ===== time embeddings (month/hour shared by both models; e_t only for gen) =====
\node[emb] (temb) at (3.9,1.7) {$e_{\text{month}}+e_{\text{hour}}$};
\node[emb] (et)   at (5.6,1.7) {$e_t$};
\draw[flow] (temb) -- (det);

% ===== lane 2: residual flow (T Euler steps) =====
\node[state] (z0) at (-4.7,-0.8) {$z^{(0)}\!\sim\!\mathcal{N}(0,\sigma^{2}\mathbf{I})$};
\node[state] (zi) at (-2.3,-0.8) {$z^{(t)}$};
\node[state] (catgen) at (0.6,-0.8)
  {$\bigl[\, z^{(t)} \,\|\, x_{n-1} \,\|\, x_{n} \,\|\, \hat{x}^{\text{det}} \,\bigr]$};
\node[model] (gen) at (3.9,-0.8) {generative model};
\node[state] (ri) at (6.4,-0.8) {$r_t$};
\node[state] (ui) at (6.4,-2.5) {$u_t=\dfrac{r_t-z^{(t)}}{s_t}$};
\node[state] (euler) at (2.6,-2.5) {$z^{(t+1)}=z^{(t)}+h_t\,u_t$};

\draw[flow] (temb) -- (gen);
\draw[flow] (et.south) -- ($(gen.north)+(0.55,0)$);
\draw[flow] (z0) -- node[lbl, above] {$t=0$} (zi);
\draw[flow] (zi) -- (catgen);
\draw[flow] (catgen) -- (gen);
\draw[flow] (gen) -- (ri);
\draw[flow] (ri) -- (ui);
\draw[flow] (ui) -- (euler);
\draw[loop] (euler.west) -| node[lbl, pos=0.25, below] {repeat $t=0,\dots,T\!-\!1$} (zi.south);

% conditioning wires into the concatenation (the det output is re-used!)
\draw[flow] (xprev.south) to[out=-90, in=90] ($(catgen.north)+(-0.35,0)$);
\draw[flow] (xcur.south)  to[out=-90, in=90] ($(catgen.north)+(0.4,0)$);
\draw[flow, rounded corners=4pt] (xdet.south) |- (1.8,0.6)
  -- ($(catgen.north)+(1.2,0)$);

% ===== composition =====
\node[op]    (plus)   at (10.9,3.0) {$\oplus$};
\node[state] (xfinal) at (12.2,3.0) {$\hat{x}_{n+1}$};
\node[state] (zT)     at (10.9,-3.7) {$\sigma_r \odot z^{(T)}$};

\draw[flow] (euler.south) |- node[lbl, pos=0.75, below] {after $T$ steps} (zT.west);
\draw[flow] (zT.north) -- (plus.south);
\draw[flow] (xdet.east) -- (plus.west);
\draw[flow] (plus) -- (xfinal);

% caption line
\node[lbl, anchor=north] at (3.7,-4.5)
  {$s_t$: noise level \quad $h_t$: Euler step size \quad $\sigma_r$: residual scale
   \quad $\hat{x}_{n+1}=\hat{x}^{\text{det}}+\sigma_r \odot z^{(T)}$};

\end{tikzpicture}%

    }
    \caption{
        Computational graph of one ArchesWeather forecast step.
        The deterministic model produces the base forecast
        \(\hat{x}^{\mathrm{det}}_{n+1}\), while the flow model generates the
        residual through \(T\) Euler steps before both components are combined.
    }
    \label{fig:arches-computational-graph}
\end{figure}

This structure will be central to our guidance method, which acts during the
residual flow while ultimately steering the resulting weather forecast.